\chapter*{Lời nói đầu}
\addcontentsline{toc}{chapter}{Lời nói đầu}

Mất cân bằng dữ liệu là một trong những thách thức dai dẳng và mang tính nền tảng của học máy ứng
dụng. Trên thực tế, rất hiếm khi ta thu thập được một bộ dữ liệu cân bằng hoàn hảo: số ca bệnh hiếm
luôn ít hơn ca bệnh phổ thông, số loài sinh vật quý hiếm luôn ít hơn loài thường gặp, và số giao
dịch gian lận bao giờ cũng là thiểu số so với giao dịch hợp lệ. Vì vậy phân bố tần suất lớp thường
có dạng \emph{đuôi-dài} (long-tailed), trong đó một số ít lớp \emph{phổ biến} (head) chiếm phần lớn
dữ liệu, còn rất nhiều lớp \emph{hiếm} (tail) chỉ có một lượng mẫu rất nhỏ. Trớ trêu là chính những
lớp hiếm này lại thường là phần quan trọng nhất đối với ứng dụng.

Trong khuôn khổ học phần \emph{Học máy và Khai phá dữ liệu} (IT3190), báo cáo trình bày kết quả
nghiên cứu của nhóm về bài toán \textbf{nhận dạng ảnh trên dữ liệu đuôi-dài} (long-tailed image
recognition). Thay vì chỉ áp dụng một thuật toán đơn lẻ, nhóm chọn cách tiếp cận như một
\emph{nghiên cứu so sánh có hệ thống}: triển khai và đánh giá một phổ rộng các phương pháp xử lý mất
cân bằng, trải dài từ những kỹ thuật huấn luyện-từ-đầu kinh điển cho tới việc thích nghi các mô hình
nền tảng (foundation model) hiện đại. Mục tiêu là trả lời câu hỏi: với các lớp đuôi chỉ có vài ảnh,
đâu là cách hiệu quả nhất để cải thiện độ chính xác, và loại tri thức bên ngoài nào giúp ích nhiều
nhất?

Toàn bộ mã nguồn, dữ liệu, log huấn luyện cũng như kết quả thực nghiệm được dẫn ra trong báo cáo đều
là sản phẩm thực tế của nhóm, chạy trên hai bộ dữ liệu CIFAR-100-LT và CUB-200-LT. Nhóm xin chân
thành cảm ơn TS.~Ngô Văn Linh đã tận tình hướng dẫn và góp ý trong suốt quá trình thực hiện.

\vspace{0.5cm}
\begin{flushright}
\textit{Nhóm sinh viên thực hiện}
\end{flushright}
